% appendix_e.tex — Extraction Table (Abbreviated)
\section{Extraction Table (Abbreviated)}\label{app:extraction}

The full extraction table (\texttt{analysis/extraction\_schema.csv}, 185
data rows + 1 header) records one row per (paper $\times$ model $\times$
dataset $\times$ metric) observation. The complete CSV is available in the
supplementary materials. Table~\ref{tab:extraction-columns} shows the
key columns; columns omitted for space include
\texttt{dataset\_variant}, \texttt{n\_series},
\texttt{series\_length\_median}, \texttt{frequency},
\texttt{fine\_tuned}, \texttt{runtime\_reported},
\texttt{gpu\_hours}, \texttt{hardware}, and \texttt{notes}.

\subsection{Column Definitions}

\begin{table}[htbp]
\centering
\caption{Key columns in the extraction table.}\label{tab:extraction-columns}
\footnotesize
\begin{tabular}{lp{8.5cm}}
\toprule
Column & Description \\
\midrule
\texttt{paper\_id}       & Unique row identifier: \texttt{\{category\}\{seq\}\_\{detail\}} \\
\texttt{authors}         & First author et al. \\
\texttt{year}            & Publication year \\
\texttt{dataset}         & Target dataset (M5, Favorita, Rohlik, GIFT-Eval, fev-bench, etc.) \\
\texttt{model\_name}     & Model as named by the paper \\
\texttt{model\_family}   & One of: foundation, deep\_learning, ml\_tree, statistical, ensemble \\
\texttt{zero\_shot}      & Yes / No --- whether the model was applied without task-specific training \\
\texttt{horizon}         & Forecast horizon (days) or ``mixed'' if aggregated \\
\texttt{eval\_method}    & rolling / fixed\_origin / competition \\
\texttt{metric\_name}    & Reported metric (\WAPE{}, \WRMSSE{}, \MASE{}, MAE, sMAPE, WQL, CRPS, win\_rate\_*) \\
\texttt{metric\_value}   & Reported value (numeric or qualitative rank) \\
\texttt{has\_covariates} & Whether the model used exogenous covariates \\
\bottomrule
\end{tabular}
\end{table}

\subsection{Row Counts by Category}

\begin{table}[htbp]
\centering
\caption{Extraction table row counts by category.}\label{tab:extraction-counts}
\begin{tabular}{llr}
\toprule
Category & Description & Rows \\
\midrule
A     & Foundation model papers            & 48 \\
B     & Benchmark papers                   & 42 \\
C     & M5 competition and retail-specific & 38 \\
D     & Deep learning baselines            & 12 \\
E     & Surveys and meta-studies           & 7 \\
F     & Cost/efficiency analysis           & 8 \\
OWN   & Our LightGBM baselines (Source~B)  & 12 \\
LOCAL & Our FM experiments (Source~B)       & 18 \\
\midrule
\textbf{Total} & & \textbf{185} \\
\bottomrule
\end{tabular}
\end{table}

\subsection{Source A vs Source B}

\begin{itemize}[nosep]
  \item \textbf{Source~A} (literature extraction): categories A--F
    (155 rows from 38 studies). These rows report accuracy numbers as
    published by the original authors.
  \item \textbf{Source~B} (gap-filling experiments): OWN + LOCAL
    (30 rows). These rows are from our own experiments on consumer
    hardware (MacBook M-series MPS) and Azure ML (E4DS\_V4), using the
    shared protocol of \S5.1.
\end{itemize}

The extraction table is frozen at the version used for the
meta-regression of \S6. Any post-submission additions would require
re-running the pooled models of \S6.7.
